\scanpage{176}
\mspage{34}

Soit $\xi_{i_0}^{\unsure{\xi|\vphantom{\xi}}}$ dans $\operatorname{Hom}(b,x_{i_0})$, ses deux images dans $\operatorname{Hom}(a,x_{i_{\unsure{0|\vphantom{0}}}})$ ne sont pas néc. égales, mais le \unsure{deviennent|sentent} dans la limite inductive, réduite à $e$, donc \deleted{qu'elles sont égales} \inserted{deviennent égales en} composant avec $x_{i_0}\to x_i$, pour $i$ assez grand, donc le composé $b\unsure{\mathbin{-}|\to}x_{i_0}\mathbin{-}x_i$

\[
a\overset{\displaystyle\longrightarrow}{\unsure{\longrightarrow|-}}b\longrightarrow x_{i_0}\mathrel{\unsure{\to|-}}x_i
\]

\deleted{\unsure{il en}} égalise la double flèche.

\uline{Corollaire} 2.16. Soit $W$ un localiseur fondamental satisfaisant l'axiome des limites $\mathrm{Loc}\,(\unsure{d^\circ|*|L^\circ})$. Si $X$ \inserted{dans Cat} est une catégorie filtrante ou cofiltrante, alors $X$ est $W$-asphérique.

Plus généralement et agréablement :

% slanted
\reversemarginpar\margin{cf $\underline{\overline{\mathrm{XII}}}$ \S\ \uline{V} : pour cette notion, que j'ai beaucoup peiné à dégager et à comprendre !}\normalmarginpar

\uline{Corollaire} 2.17. Soit $W$ comme dans 2.16. Alors \deleted{\unsure{toute}} catégorie filtrante \uline{\unsure{totalement} $W$-asphérique}, i-e pour \deleted{tout, foncteur} \unsure{flèche} $f:X\longrightarrow Y$ dans Cat, \deleted{la catégorie \unsure{$f_!(X/Y)$}} l'objet $f_!(e_X)$ de $e_Y$ \deleted{est $W$-asphérique,} (défini aussi comme \deleted{l'espace} $y\longmapsto\pi_0({}_{y}\backslash X)$, et intervenant dans la factorisation canonique de $f$ en \deleted{\unsure{épi}}, en \unsure{p.} étale et \unsure{en|les} \unsure{cofibration|cofibres} \deleted{\unsure{de|des}} \unsure{fibr. à|fes b} \inserted{$0$-connexes})

% archive sheet number: 176
